% TOP_PROBLEM: {"id": 30006728, "title": "Grothendiecks Section Conjecture in Anabelian Geometry", "category_id": 5, "category": {"id": 5, "name": "algebraic_geometry", "display_name": "Algebraic Geometry", "description": "Geometric objects defined by polynomial equations.", "slug": "algebraic-geometry", "order_index": 5, "created_at": "2026-07-31T15:26:25.671Z"}, "status": "open"}
% FIRST_POSTED: 2026-10-01
% ATTEMPT_STATUS: unresolved
% !TeX program = lualatex
% Definition and sources only; this file is not an LLM solution attempt.
\documentclass[11pt]{article}
\usepackage[margin=25mm]{geometry}
\usepackage{fontspec}
\setmainfont{Latin Modern Roman}
\usepackage{amsmath,amssymb,mathtools}
\usepackage{unicode-math}
\setmathfont{Latin Modern Math}
\usepackage{xurl,hyperref}
\hypersetup{colorlinks=true,urlcolor=blue}
\setcounter{secnumdepth}{0}
\setlength{\parindent}{0pt}
\setlength{\parskip}{5pt}
\setlength{\emergencystretch}{3em}
\begin{document}
\section*{\texorpdfstring{Grothendiecks Section Conjecture in Anabelian Geometry}{Grothendiecks Section Conjecture in Anabelian Geometry}}\label{30006728}
\subsection{Definitions and mathematical statement}
For a smooth projective geometrically connected curve $C/K$ of genus at least two, choose a geometric base point and consider
\[
 1\longrightarrow\pi_1(C_{\bar K})\longrightarrow\pi_1(C)
 \longrightarrow G_K\longrightarrow1.
\]
A section is a continuous homomorphism $s:G_K\to\pi_1(C)$ whose composition with the projection is the identity. Sections are identified under conjugation by $\pi_1(C_{\bar K})$. A rational point supplies such a conjugacy class through functoriality of the étale fundamental group. The conjecture says that the map from $C(K)$ to these section classes is bijective.
\par\smallskip\textit{Formulation and concept references:} \href{https://doi.org/10.1017/CBO9780511758874.005}{[S1]}.

\subsection{Short English statement}
Do rational points on a hyperbolic number-field curve correspond exactly to the sections of its arithmetic fundamental-group sequence?
\subsection{Research attempt: nonabelian cocycles and a concrete lifting obstruction}
\textbf{Status.} We give an exact description of all sections conditional on one existing section, and exhibit why abelianized tests cannot be sufficient in general. These are reductions of the section problem, not a proof that all sections of a curve are point-theoretic. The primary arithmetic treatment [R1] concerns particular birational section results with additional hypotheses, which are not silently removed here.

\subsubsection{1. Sections relative to a chosen section}
Write the exact sequence as $1\to N\to E\to G\to1$ and suppose a continuous section $s_0$ is given. For instance, a known $K$-point supplies one in the target setting. Define the continuous action $g\cdot n=s_0(g)ns_0(g)^{-1}$. Every section has a unique expression
\[s(g)=c(g)s_0(g),\qquad c:G\to N\text{ continuous}.\]
Multiplying two such expressions shows that $s$ is a homomorphism exactly when
\[c(gh)=c(g)(g\cdot c(h)).\tag{1}\]
Conversely, any continuous function satisfying (1) defines a continuous section. Conjugating $s$ by $n\in N$ changes the cocycle to
\[c'(g)=n c(g)(g\cdot n^{-1}).\tag{2}\]
Hence section classes are exactly the nonabelian continuous cohomology pointed set $H^1(G,N)$ for this action. This is a bijection obtained from the equations, not merely an injection into a cohomological invariant.

For $A=N^{\mathrm{ab}}$, passage to the abelianization sends each such class to $H^1(G,A)$. In additive notation, (1) becomes $a(gh)=a(g)+g\cdot a(h)$ and (2) differs by a coboundary. Thus failure of an abelianized necessary condition can obstruct a proposed section. Satisfaction of that condition need not produce a nonabelian lift.

\subsubsection{2. An explicit finite obstruction to lifting abelian data}
For an odd prime $p$, let $H$ be the group of upper unitriangular $3\times3$ matrices over $\mathbb F_p$, written as triples $(a,b,c)$ with multiplication
\[(a,b,c)(a',b',c')=(a+a',b+b',c+c'+ab').\]
The map $H\to\mathbb F_p^2$, $(a,b,c)\mapsto(a,b)$, has central kernel. Suppose the identity homomorphism of $\mathbb F_p^2$ lifted to a homomorphism $\ell:\mathbb F_p^2\to H$. The two coordinate basis vectors commute, so their lifts would commute. But any lifts have forms $(1,0,c)$ and $(0,1,d)$, and their commutator is $(0,0,1)\ne1$. Thus no such lift exists.

The quotient is abelian and the proposed quotient-level map is as simple as possible, yet a central commutator prevents a lift. All groups are finite, so continuity creates no extra issue in this example. This is an elementary model of the nonlinear obstruction lost by abelianization; it is not a counterexample to Grothendieck's conjecture for arithmetic fundamental groups.

\subsubsection{3. What the arithmetic problem becomes}
If $C(K)$ is nonempty and $s_0$ comes from a chosen point, the conjectured surjectivity says every class in the explicitly defined $H^1(G_K,\pi_1(C_{\bar K}))$ is represented by some rational point. This requires identifying a very special geometric subset of a nonabelian cohomology set. Merely showing that a class has plausible images in abelian or finite quotients does not identify that subset. Compatibility and lifting across quotients can carry further obstructions.

If $C(K)$ is empty, no point supplies $s_0$. The semidirect-product parametrization then cannot be invoked without first proving that some section exists. In that case the conjecture predicts that none exists. Assuming a splitting as a harmless choice of notation would already assume the central disputed fact.

\subsubsection{4. Audit and remaining gap}
The conjugation equivalence is by the geometric subgroup $N$, exactly as in the target, not by arbitrary elements of $E$. The cocycles and sections are continuous, as required for profinite groups. The finite lifting calculation was checked directly from the multiplication formula. What remains is the arithmetic assertion that all continuous sections are induced by points and that distinct rational points give distinct classes; the general cocycle parametrization proves neither.

\subsection{Sources}
\begin{itemize}
\item[S1]A. Grothendieck, Brief an G. Faltings, LMS Lecture Notes 242:49-58, 1997.
\url{https://doi.org/10.1017/CBO9780511758874.005}.
\textit{Formulation source.}
\item[S2]Polylogarithmic Chabauty--Kim loci over number fields.
\url{https://arxiv.org/abs/2608.20615}.
\textit{Further reference; not a proof endorsement.}
\item[S3]On the section conjecture over fields of finite type.
\url{https://onlinelibrary.wiley.com/doi/full/10.1002/mana.70049}.
\textit{Further reference; not a proof endorsement.}
\item[R1] Jakob Stix, On the birational section conjecture with local conditions (primary research preprint).
\url{https://arxiv.org/abs/1203.3236}.
\end{itemize}
\end{document}
